% How to get the course repository: clone it with git, or download a ZIP.
% Build: pdflatex getting-started.tex
\documentclass[11pt,a4paper]{article}
\usepackage[margin=2.2cm]{geometry}
\usepackage[T1]{fontenc}
\usepackage{lmodern}
\usepackage[skins]{tcolorbox}
\usepackage{fontawesome5}
\usepackage{xcolor}
\usepackage{enumitem}
\usepackage[hidelinks]{hyperref}
\usepackage{titlesec}

\definecolor{ink}{HTML}{1F2A37}
\definecolor{accent}{HTML}{D9622B}
\definecolor{link}{HTML}{1F5FBF}
\definecolor{term}{HTML}{F3F4F6}

\titleformat{\section}{\large\bfseries\color{ink}}{\color{accent}\thesection}{0.6em}{}
\setlength{\parindent}{0pt}
\setlength{\parskip}{6pt}
\pagestyle{empty}

% a terminal: commands to type, one per line
\newtcolorbox{terminal}{enhanced, colback=term, colframe=ink!30, boxrule=0.6pt, arc=3pt,
  left=8pt, right=8pt, top=4pt, bottom=4pt, fontupper=\ttfamily\small}
\newtcolorbox{tip}{enhanced, colback=accent!8, colframe=accent, boxrule=0.6pt, arc=3pt,
  left=8pt, right=8pt, top=4pt, bottom=4pt, before upper=\raggedright}
\newcommand{\link}[1]{\href{#1}{\textcolor{link}{\texttt{#1}}}}

\begin{document}

{\LARGE\bfseries\color{ink} Getting the course files}\par\vspace{2pt}
{\large Cryptography: stream ciphers, RSA, Diffie-Hellman, elliptic curves}\par\vspace{6pt}
All slides, tasks and starter files are in one place on GitHub:\par
\begin{center}\Large\faGithub\ \link{https://github.com/KsaweryM/Cryptography-Course}\end{center}
There are two ways to get them. \textbf{Way A} (git) is better: one command gets the updates later.
\textbf{Way B} (ZIP) needs nothing installed.

\section{Way A: clone with git}

\textbf{Step 1. Check that git is installed.} Open a terminal
(Windows: \emph{Git Bash} or \emph{PowerShell}; macOS: \emph{Terminal}; Linux: any terminal) and type:
\begin{terminal}
git -{}-version
\end{terminal}
If you see a version number, for example \texttt{git version 2.43.0}, go to step 2.
If not, install git:
\begin{itemize}[nosep, leftmargin=1.4em]
  \item \textbf{Windows:} download it from \link{https://git-scm.com/download/win} and click through the installer (the default options are fine).
  \item \textbf{macOS:} type \texttt{xcode-select -{}-install} and confirm.
  \item \textbf{Linux (Ubuntu, Debian):} type \texttt{sudo apt install git}.
\end{itemize}
Then close the terminal, open a new one and check again.

\textbf{Step 2. Go to the folder where you keep your work}, for example your Documents:
\begin{terminal}
cd Documents
\end{terminal}

\textbf{Step 3. Clone the repository.} This copies everything into a new folder \texttt{Cryptography-Course}:
\begin{terminal}
git clone https://github.com/KsaweryM/Cryptography-Course.git\\
cd Cryptography-Course
\end{terminal}
No GitHub account and no password are needed: the repository is public.

\textbf{Step 4. Later: get the new files.} Before each lecture, inside the folder \texttt{Cryptography-Course}:
\begin{terminal}
git pull
\end{terminal}

\begin{tip}
\faLightbulb\ \textbf{If \texttt{git pull} complains} that your changes would be overwritten:
you edited a file that was also updated. Keep your work with \texttt{git stash}, then \texttt{git pull},
then \texttt{git stash pop}. Or simply copy your solution somewhere else first.
\end{tip}

\section{Way B: download a ZIP}

\begin{enumerate}[nosep, leftmargin=1.6em]
  \item Open \link{https://github.com/KsaweryM/Cryptography-Course} in a browser.
  \item Click the green button \textbf{\textless\textgreater\ Code}, then \textbf{Download ZIP}.
  \item Unpack the ZIP file (right click, \emph{Extract all} on Windows).
\end{enumerate}
To get updates, download the ZIP again. Keep your own solutions in a separate folder, so a new ZIP never overwrites them.

\section{What is inside, and how to run a task}

\begin{itemize}[nosep, leftmargin=1.4em]
  \item \texttt{Lecture1} \dots\ \texttt{Lecture5}: one folder per lecture, with the slides as a PDF.
  \item \texttt{LectureN/TaskK}: one folder per task. Start with its \texttt{README.md}: it tells you the goal,
        the steps and the expected output. On GitHub the README is shown with its picture under the file list.
\end{itemize}

The tasks need only \textbf{Python 3.8 or newer}, nothing else to install. Check it, then run a task, for example Task 1 of Lecture 1:
\begin{terminal}
python3 -{}-version\\
cd Lecture1/Task1\\
python3 decrypt.py
\end{terminal}
On Windows the command is often \texttt{py} instead of \texttt{python3} (\texttt{py -{}-version}, \texttt{py decrypt.py}).
No Python? Get it from \link{https://www.python.org/downloads/}; on Windows tick \emph{Add python.exe to PATH} in the installer.

\begin{tip}
\faDesktop\ \textbf{Interactive simulators.} The slides also link to pages that show each algorithm step by step,
in the browser, nothing to install: \link{https://ksawerym.github.io/Cryptography-simulators/}
\end{tip}

\end{document}
